\uuid{vdkX}
\exo7id{940}
\titre{exo7 940}
\auteur{cousquer}
\organisation{exo7}
\datecreate{2003-10-01}
\isIndication{false}
\isCorrection{false}
\chapitre{Application linéaire}
\sousChapitre{Image et noyau, théorème du rang}
\cpge{mpsi-mp2i.applications-lineaires.a, pcsi.applications-lineaires.a, ptsi.applications-lineaires.a, mpsi-mp2i.espaces-vectoriels.c, pcsi.espaces-vectoriels.c, ptsi.espaces-vectoriels.c, mpsi-mp2i.applications-lineaires.d, pcsi.applications-lineaires.d, ptsi.applications-lineaires.d, bcpst2.espaces-vectoriels.a, bcpst2.applications-lineaires-matrices.a}
\module{Algèbre}
\niveau{L1}
\difficulte{}

\contenu{
\texte{
Soit $E$ un espace vectoriel, et $u$ une application linéaire de~$E$ 
dans~$E$. Dire si les
propriétés suivantes sont vraies ou fausses~:
}
\begin{enumerate}
    \item \question{Si $e_1,e_2,\ldots,e_p$ est libre, il en est de même de
$u(e_1),u(e_2),\ldots,u(e_p)$.}
    \item \question{Si $u(e_1),u(e_2),\ldots,u(e_p)$ est libre, il en est de même de
$e_1,e_2,\ldots,e_p$.}
    \item \question{Si $e_1,e_2,\ldots,e_p$ est génératrice, il en est de même de
$u(e_1),u(e_2),\ldots,u(e_p)$.}
    \item \question{Si $u(e_1),u(e_2),\ldots,u(e_p)$ est génératrice, il en est de même de
$e_1,e_2,\ldots,e_p$.}
    \item \question{Si $u(e_1),u(e_2),\ldots,u(e_p)$ est une base de $\Im u$, alors
$e_1,e_2,\ldots,e_p$ est une base d'un sous-espace vectoriel supplémentaire
de $\mathrm{Ker} u$.}
\end{enumerate}
}
